% GPT-6 Astra Pro; effort not exposed; /root/replacement; 22 November 2026.
\section[HD-F1-R32: the quotient target is simple]{HD-F1-R32: the actual common-quotient target is ordinarily simple}
\label{sec:followup-rm-common-quotient-simple-target}

Every actual object in the common-quotient branch
of R21--R31 has
\[
                         \operatorname{End}_X(T)=\mathbf C.
 \tag{F1R32.1}
\]
This is ordinary simplicity of the target
coherent sheaf. In particular this branch
contains no target that is an extension of
one rank-one TF sheaf by itself. More
generally, every actual rank-one TF
filtration of $T$ has zero reverse Hom
between its quotient and subobject, as
proved below. The original coupling,
point action and surface defects are
retained, and their remaining existence
problem is not decided here.

\paragraph{1. Ordinary endomorphisms act on the original invariant columns.}
Keep the actual maps of (F1R31.1)--(F1R31.3):
\[
 p:A\twoheadrightarrow Q,\quad l:Q\hookrightarrow L=Q^{**},
 \quad u_1,u_2:Q\to T,\quad
 j:T\hookrightarrow T^{**}=L^{\oplus2},
 \qquad j_a u_b=\delta_{ab}l.
 \tag{F1R32.2}
\]
Double dual induces the actual injection
\[
 \operatorname{Hom}_X(Q,T)
       \lhook\joinrel\longrightarrow
       \operatorname{Hom}_X(L,L^{\oplus2})=\mathbf C^2.
 \tag{F1R32.3}
\]
A map in its kernel is zero on the
common locally free open and hence
zero globally, since $T$ is torsion
free. The two original invariant maps
$u_1,u_2$ go to the standard basis.
They therefore span the entire
ordinary Hom space as well:
\[
        \operatorname{Hom}_X(Q,T)
             =\mathbf C u_1\oplus\mathbf C u_2
             =\operatorname{Hom}_G(Q,T).
 \tag{F1R32.4}
\]
Thus postcomposition by an arbitrary
ordinary endomorphism of $T$ remains
in this particular invariant column
space. This has been proved from the
actual double-dual map, without
assuming that endomorphism is invariant.

Likewise double dual embeds the
ordinary endomorphism algebra into
the constant hull matrices:
\[
 \operatorname{End}_X(T)\lhook\joinrel\longrightarrow
                  \operatorname{End}_X(L^{\oplus2})
                        =\operatorname{Mat}_2(\mathbf C).
 \tag{F1R32.5}
\]
For $n\in\operatorname{End}_X(T)$
let $M_n$ be this matrix. Naturality
of $j$ and (F1R32.2) give
\[
                         n u_w=u_{M_nw}
                             \quad(w\in\mathbf C^2).
 \tag{F1R32.6}
\]
Both sides have hull column $M_nw$;
the injection (F1R32.3) proves the
equality of the actual maps.

For the same original connecting
class $\delta:T\to A[2]$, put
$\beta_Q=p[2]\delta$ and define
\[
 \begin{aligned}
 \mu:\mathbf C^2&\longrightarrow\operatorname{Ext}^2_X(T,T),
                         &w&\longmapsto u_w[2]\beta_Q,\\
 \rho=\sigma_T^2\mu:&\mathbf C^2\lhook\joinrel\longrightarrow HH_{-2}(X).
 \end{aligned}
 \tag{F1R32.7}
\]
The ordinary trace in the second
line includes all Hodge components.
Its injectivity, proved in paragraph~2
of R31, follows directly from the
rank-two whole characteristic image
of R26, the actual target Yoneda
membership, and the fact that its
possible coefficient space has
dimension two. The same argument
gives a surjective linear map
\[
          v:K\twoheadrightarrow\mathbf C^2,
                    \qquad\operatorname{ev}_T^2(k)=\mu(v(k)),
 \tag{F1R32.8}
\]
where $K$ is the original balanced
degree-two Hochschild kernel.
The maps in (F1R32.7) have invariant
image, but their ordinary full trace
is the one used here.

\paragraph{2. Ordinary indecomposability reduces the matrix to a scalar and a square-zero map.}
The exact ordinary scope of R24,
paragraph~4, is used: $T$ has no
ordinary rank-one retract, whether
or not that retract preserves $G$.
Its proof restricts the actual
endpoint composition to such a
retract; the ordinary space
$\operatorname{Hom}_X(Q,F)$ has
dimension at most one, and the
resulting rank-one characteristic
action forces the prohibited
effective hull cycle $q(A_1+A_2)$.
Since every nontrivial decomposition
of a rank-two TF sheaf has rank-one
TF summands, $T$ is ordinarily
indecomposable.

If an actual hull matrix in
(F1R32.5) had two distinct eigenvalues,
the two linear polynomials in the
original endomorphism giving its
spectral projectors would be actual
nonzero complementary idempotents
of $T$. Their images give an ordinary
decomposition, contrary to that
proved scope of R24. Hence every
endomorphism is $\lambda 1_T+n$
with
\[
                              n^2=0.
 \tag{F1R32.9}
\]
Indeed its constant two-by-two matrix
has a single eigenvalue; Cayley--Hamilton
gives $(M-\lambda I)^2=0$, and
(F1R32.5) is injective.

\paragraph{3. The square-zero image factorization kills its full trace.}
Let $I=\operatorname{im}n$ be the
actual coherent image and factor
\[
                T\xrightarrow{a}I\xrightarrow{i}T,
                  \qquad a\text{ surjective},\quad
                  i\text{ injective},\quad n=ia.
 \tag{F1R32.10}
\]
The equality $n^2=i(ai)a=0$,
together with these actual epimorphism
and monomorphism maps, implies
\[
                                  ai=0.
 \tag{F1R32.11}
\]
All these sheaves are perfect on
the smooth fourfold. Let
$\operatorname{AC}_M=\exp(-\operatorname{At}_M)$
denote the retained Buchweitz--Flenner
Atiyah--Chern transformation.
Naturality and degree-zero trace
cyclicity give component by component
\[
 \begin{split}
 \sigma_T^0(n)
   &=\operatorname{Tr}_T\bigl((ia\otimes1)\operatorname{AC}_T\bigr)\\
   &=\operatorname{Tr}_T\bigl((i\otimes1)\operatorname{AC}_I a\bigr)\\
   &=\operatorname{Tr}_I\bigl(\operatorname{AC}_I ai\bigr)=0.
 \end{split}
 \tag{F1R32.12}
\]
The two moved morphisms $i,a$ have
degree zero, so there is no odd
interchange sign. This proves
vanishing of the full trace, not
only of the usual generic-rank trace.
It uses the actual image factorization
rather than a claim about traces of
arbitrary abstract nilpotents.

For precision about the next action,
let $\sigma_M^H$ be the Hochschild-valued
trace. The actual decomposition
isomorphism $\Phi$ gives
$\sigma_M=\Phi\sigma_M^H$ by
Buchweitz--Flenner Proposition~6.2.1.
Thus (F1R32.12) also gives
$\sigma_T^H(n)=0$.
Their Proposition~6.3.1 is the actual
module identity
\[
       \sigma_T^H(\operatorname{ev}_T^2(k)n)
                             =k\cdot\sigma_T^H(n)=0.
 \tag{F1R32.13}
\]
The graded-centre action commutes
with the degree-zero map $n$.
Applying the same $\Phi$ therefore
gives the ordinary full-trace equality
\[
                          \sigma_T^2(n[2]\operatorname{ev}_T^2(k))=0.
 \tag{F1R32.14}
\]
This use of Hochschild module linearity
does not assume that a separate HKR
identification is an algebra map on
arbitrary cohomology classes.

Substitute the original maps
(F1R32.6)--(F1R32.8) into
(F1R32.14). The result is
\[
          0=\sigma_T^2\bigl(u_{M_n v(k)}[2]\beta_Q\bigr)
                              =\rho(M_n v(k))
                               \quad(k\in K).
 \tag{F1R32.15}
\]
Injectivity of $\rho$ and surjectivity
of $v$ force $M_n=0$. Injectivity
of the actual hull map (F1R32.5)
then forces $n=0$. Together with
paragraph~2 this proves (F1R32.1).

\paragraph{4. The excluded actual target extensions.}
Consider any actual short exact
sequence of rank-one TF sheaves
\[
                0\longrightarrow F\xrightarrow{b}T
                          \xrightarrow{c}R\longrightarrow0.
 \tag{F1R32.16}
\]
For a nonzero ordinary morphism
$f:R\to F$, the composite
$bfc$ is a nonzero endomorphism of
$T$, because $c$ is surjective and
$b$ is injective. It satisfies
\[
                        (bfc)^2=bf(cb)fc=0.
 \tag{F1R32.17}
\]
This contradicts (F1R32.1).
Consequently every sequence
(F1R32.16) satisfies
\[
                               \operatorname{Hom}_X(R,F)=0.
 \tag{F1R32.18}
\]
In particular $F\simeq R$ is
excluded, whether the extension is
split or nonsplit, and whether its
filtration is equivariant or ordinary.

There is a useful ideal-sheaf version
for the slope-zero saturated filtrations
provided by the common line hull.
For such an $F$, a nonzero component
of $F^{**}\to L^{\oplus2}$ is a
map of slope-zero lines. Its effective
zero divisor has ample degree zero,
so it is an isomorphism $F^{**}\simeq L$.
The determinant isomorphism from
(F1R32.16) gives $R^{**}\simeq L$.
Use these actual hull maps to write
\[
                         F=\mathcal I_1L,\qquad
                         R=\mathcal I_2L.
 \tag{F1R32.19}
\]
Double dual injects $\operatorname{Hom}(R,F)$
into $\operatorname{End}(L)=\mathbf C$.
A nonzero such map is a nonzero
scalar on the common big open;
division by that scalar extends
the identity on the hull and is
exactly an inclusion
$\mathcal I_2\subset\mathcal I_1$.
Conversely that inclusion defines
an actual nonzero map $R\to F$.
Thus (F1R32.18) excludes precisely
that inclusion direction in every
such actual target filtration. In
particular the two ideals cannot be
equal. No assertion of the reverse
inclusion or of the existence of a
remaining extension is made.

\paragraph{Framework and checking scope.}
The ordinary retract argument is the
already proved R24 paragraph~4, whose
full proof was reread at its ordinary
scope. The original hull and column
maps are R22 and R31; the injective
trace map and surjective coefficient
map are (F1R31.8)--(F1R31.12).
The new obstruction uses the actual
image maps (F1R32.10), their vanishing
reverse composition, and their full
cyclic trace.

The published framework is
Buchweitz--Flenner, \emph{The global
decomposition theorem for Hochschild
(co-)homology of singular spaces via
the Atiyah--Chern character},
Adv.~Math.~217 (2008), 243--281,
DOI 10.1016/j.aim.2007.06.013.
The complete proofs of Propositions
6.2.1 and 6.3.1 were freshly read
in the primary author version,
pp.~29--30, including the exact
trace diagram and graded-centre
module calculation. Naturality of
the Atiyah--Chern transformation and
the degree-zero trace sign are the
retained published framework, with
the exact needed compositions
written in (F1R32.12)--(F1R32.15).
This excludes additional actual
target extensions; the positive-pure
surface cancellation and coupled
defect extension remaining after
R31 still require construction or
another obstruction.
